% ========================================================
% 0abstract.tex
% ========================================================
\begin{abstract}
High-density transit venues---airports, railway stations, and bus terminals---pose
unique challenges for WiFi resource management: massive concurrent users,
schedule-driven load surges, multi-class quality-of-service (\mbox{QoS}) requirements,
and the need for coordination across dozens of access points (APs) without
sharing raw user data. Existing approaches are either reactive (failing to
anticipate load spikes) or monolithic (lacking schedule awareness and
fairness guarantees).

We propose \textbf{STAG-FedSAC}: a Schedule-Conditioned Spatio-Temporal Graph
Attention Transformer coupled with a Hierarchical Federated Soft Actor-Critic
agent that jointly optimises load prediction, proactive resource allocation,
and fair bandwidth distribution across APs. Our framework introduces four
novel contributions: \emph{(i)} ST-GCAT, the first graph cross-attention
Transformer that fuses transit schedule features with spatio-temporal AP load
patterns to predict congestion up to 30 minutes ahead; \emph{(ii)} a
graph-embedded SAC state that injects structured ST-GCAT embeddings directly
into the deep reinforcement learning (DRL) policy, providing spatial context
beyond scalar load estimates; \emph{(iii)} a bidirectional joint training
mechanism in which policy gradients back-propagate through the predictor
($\mathcal{L}_\mathrm{pred} = \mathrm{MSE} + \eta \mathcal{L}_\mathrm{actor}$),
creating a tighter prediction--control coupling; and \emph{(iv)} HierFed-KD,
a two-level (zone/global) federated aggregation protocol enhanced by
knowledge distillation that achieves a $11.5\times$ communication reduction
relative to raw-parameter sharing while preserving personalised per-AP policies.

Evaluated on a 30-AP simulated transit terminal over 2000 training episodes,
STAG-FedSAC achieves an episode reward of $0.763$, a Jain Fairness Index (JFI)
of $0.921$, a mean per-user throughput of $74.8$~Mbps, and an average latency of
$32.4$~ms---improvements of $24.5\%$, $9.5\%$, $18.0\%$, and $27.5\%$,
respectively, over the strongest federated baseline (FedDDPG). Load prediction
attains RMSE $= 0.041$ and Surge-RMSE $= 0.067$, satisfying the target
thresholds with and without schedule-driven surge events. Ablation studies
confirm that each novel component contributes measurably to overall performance.

\end{abstract}
